La trajectòria de la Terra al voltant del Sol és una el·lipse; aquest fet fa que la distància des de la Terra al Sol no sigui la mateixa en totes les èpoques de l'any. El periheli, la distància més curta entre la Terra i el Sol, és de $1,471\times 10^{8}\ \mathrm{km}$. La Terra passa pel periheli durant els primers dies del mes de gener de cada any. La velocitat de la Terra al periheli és de $30,75\ \mathrm{km/s}$. L'afeli és la posició més allunyada del Sol. Quan la Terra es troba a l'afeli, la seva velocitat orbital és de $28,76\ \mathrm{km/s}$.

\begin{apartat}{1,25}
Dibuixeu una òrbita clarament el·líptica (no cal que sigui l'òrbita real) on s'indiqui la posició del Sol i la de la Terra un dia d'hivern de l'hemisferi nord. Utilitzant arguments basats en l'energia, justifiqueu per què la velocitat de la Terra és mínima a l'afeli. Quina és la distància de la Terra al Sol a l'afeli?
\end{apartat}

\begin{apartat}{1,25}
Quina intensitat de camp gravitatori genera el Sol a la seva superfície? Quin és el pes d'una massa de $10,0\ \mathrm{kg}$ a la superfície del Sol?
\end{apartat}

\dades{%
  $G = 6,67\times 10^{-11}\ \mathrm{N\,m^2\,kg^{-2}}$.\\
  $M_\mathrm{Terra} = 5,97\times 10^{24}\ \mathrm{kg}$.\\
  $M_\mathrm{Sol} = 1,99\times 10^{30}\ \mathrm{kg}$.\\
  $R_\mathrm{Sol} = 6,96\times 10^{5}\ \mathrm{km}$.}
